\documentclass[11pt]{article}
\usepackage[margin=1in]{geometry}
\usepackage{amsmath,amssymb,amsthm,hyperref}
\newtheorem{lemma}{Lemma}
\title{Counterexamples to both squared operator inequalities\newline in MathDB 329636}
\author{Principia Math in Progress}
\date{12 September 2026}
\begin{document}
\maketitle

\section*{Task 1. Resolve the two proposed squared inequalities}
For positive definite matrices, write
$A\#B=A^{1/2}(A^{-1/2}BA^{-1/2})^{1/2}A^{1/2}$.
Conjecture 2.16 of Moslehian and Fu \cite{MF} proposes that every unital
positive linear map $\Phi$, with
$m_1^2I\leq A\leq M_1^2I$ and $m_2^2I\leq B\leq M_2^2I$,
satisfies both
\begin{align}
 (\Phi(A)\#\Phi(B))^2&\leq
 \frac{(M_1M_2+m_1m_2)^2}{4M_1M_2m_1m_2}\,\Phi(A\#B)^2,
 \label{PS}\\
 \left(\frac{M_2m_2}{M_1m_1}\Phi(A)+\Phi(B)\right)^2
 &\leq\left(\frac{M_2}{m_1}+\frac{m_2}{M_1}\right)^2\Phi(A\#B)^2.
 \label{DM}
\end{align}
\textbf{Both assertions are false.} We give separate exact counterexamples.
All inputs are positive definite and both maps are unital and completely positive.
The order is the usual positive semidefinite order.

\subsection*{A geometric-mean computation}
\begin{lemma}
Let $C>0$, $d>0$, $x\neq0$, and
$P_\pm=\begin{pmatrix}C&\pm x\\\pm x^*&d\end{pmatrix}>0$.
Put $t=x^*C^{-1}x$ and $\rho=\sqrt{1-t/d}$. Then
\[
 P_+\#P_-=\operatorname{diag}(Y,d\rho),\qquad
 Y=C-\frac{1-\rho}{t}xx^*.
\]
\end{lemma}
\begin{proof}
Congruence by $T^{-1}=\operatorname{diag}(C^{-1/2},d^{-1/2})$
gives matrices $W_\pm=\begin{pmatrix}I&\pm w\\\pm w^*&1\end{pmatrix}$,
where $w=C^{-1/2}x/\sqrt d$. The Schur complement gives $0<\|w\|^2<1$.
The two matrices commute, and their product is
$\operatorname{diag}(I-ww^*,1-w^*w)$. Thus their geometric mean is
\[
 \operatorname{diag}\left(I-\frac{1-\rho}{\|w\|^2}ww^*,\rho\right).
\]
Congruence back by $T$ proves the formula. Congruence invariance follows
directly from the characterization of $A\#B$ as the unique positive definite
$X$ satisfying $XA^{-1}X=B$: conjugating this equation by any invertible
matrix preserves the characterization. Uniqueness follows by taking the
positive square root of $A^{-1/2}BA^{-1/2}$.
\end{proof}

\subsection*{Counterexample to the squared P\'olya--Szeg\H{o} assertion}
Take $m_1=m_2=1$, $M_1=M_2=32$, and set
\[
 a=\frac{1025}{2},\quad b=\frac{1023}{2},\quad
 A=\begin{pmatrix}a&0&b\\0&1&0\\b&0&a\end{pmatrix},\quad
 B=\begin{pmatrix}a&0&-b\\0&1&0\\-b&0&a\end{pmatrix}.
\]
Both spectra are $\{1,1,1024\}$, so all spectral bounds hold.
Let $U$ interchange the first two coordinates and fix the third, and define
\[
 \Phi(X)=\frac14 X+\frac34 UXU^*.
\]
This convex combination of unitary conjugations is unital and completely
positive. Since $A$ and $B$ commute,
\[
 A\#B=\operatorname{diag}(32,1,32),\qquad
 Q:=\Phi(A\#B)=\operatorname{diag}(35/4,97/4,32).
\]
The images $\Phi(A),\Phi(B)$ are the lemma's $P_+,P_-$ with
\[
 C=\frac18\operatorname{diag}(1031,3077),\qquad
 x=\frac{1023}{8}\binom{1}{3},\qquad d=\frac{1025}{2}.
\]
Direct rational arithmetic gives
\[
 \rho^2=\frac{18968594}{3251696675},\qquad
 Y=C-c\begin{pmatrix}1&3\\3&9\end{pmatrix},\qquad
 c=\frac{3172387}{98848}(1-\rho).
\]
Comparing squares shows $7/100<\rho<8/100$. Consequently
\[
 29<\frac{72964901}{2471200}<c
 <\frac{295031991}{9884800}<30.
\]
For $z=(7,-1,0)^T$, the lemma therefore yields
\[
 \|(\Phi(A)\#\Phi(B))z\|^2
 =\left(\frac{7217}{8}-4c\right)^2
  +\left(\frac{3077}{8}+12c\right)^2
 >782^2+732^2=1147348.
\]
But the proposed right side of \eqref{PS}, evaluated on $z$, equals
\[
 \left(\frac{1025}{64}\right)^2\|Qz\|^2
 =\left(\frac{1025}{64}\right)^2
   \left[\left(\frac{245}{4}\right)^2+
         \left(\frac{97}{4}\right)^2\right]
 =\frac{36474548125}{32768}<1147348.
\]
For a Hermitian matrix $H$, $z^*H^2z=\|Hz\|^2$.
Thus the quadratic form of the proposed right side minus the left side
is negative. This disproves \eqref{PS}.

\newpage
\subsection*{Counterexample to the squared Diaz--Metcalf assertion}
Now use $A=\operatorname{diag}(1,9/4,4)$, $B=I_3$,
$m_1=1$, $M_1=2$, and $m_2=M_2=1$. Define
\[
 V=\frac13\begin{pmatrix}2&1\\1&2\\2&-2\end{pmatrix},\qquad
 \Psi(X)=V^*XV.
\]
Since $V^*V=I_2$, this is a unital completely positive map from
$\mathbb M_3$ to $\mathbb M_2$. Put $R=\Psi(A\#B)$ and
$S=\frac12\Psi(A)+I_2$. Exact multiplication gives
\[
 R=\begin{pmatrix}3/2&-1/3\\-1/3&5/3\end{pmatrix},\qquad
 S=\begin{pmatrix}161/72&-19/36\\-19/36&22/9\end{pmatrix},
\]
and
\[
 D:=\frac94R^2-S^2
 =\frac1{5184}\begin{pmatrix}175&494\\494&1276\end{pmatrix}.
\]
For $v=(3,-1)^T$,
\[
 v^*Dv=-\frac{113}{5184}<0
 \qquad\left(\text{also }\det D=-\frac1{1296}\right).
\]
This is precisely the failure of \eqref{DM}. If a map on one matrix algebra
is preferred, replace $\Psi$ by
$\widehat\Psi(X)=\operatorname{diag}(V^*XV,X_{22})$ from
$\mathbb M_3$ to itself. It is unital and completely positive, and the same
negative quadratic form occurs on $(3,-1,0)^T$.

Equality of the two bounds for $B$ is explicitly permitted in the conjecture.
The failure also persists with distinct bounds: keep $B=I_3$, and take
$M_2=10001/10000$, $m_2=10000/10001$. Their product remains one, and
the new quadratic form is exactly
\[
 v^*\left[\left(M_2+\frac{m_2}{2}\right)^2R^2-S^2\right]v
 =-\frac{55824935636320127}{3240648032400000000}<0.
\]
Hence both individual assertions in the source conjecture have been refuted.

\section*{Task 2. Determine the optimal replacement constants}
Find the smallest dimension-independent constants in \eqref{PS} and
\eqref{DM} for each prescribed set of spectral bounds, and compare them
with the weaker constants already established in the source.

\section*{Task 3. Identify maps for which the proposed bounds survive}
Characterize structural conditions on $\Phi$ and on the inputs that permit
the two unsquared inequalities to be squared with their original constants.
Include commutativity of the relevant images and multiplicativity of the map.

\section*{Task 4. Determine the smallest counterexample dimensions}
Decide whether the first assertion can fail for maps from $\mathbb M_2$
to itself, and classify the minimal input and output dimensions required
for each failure.

\section*{Task 5. Study the corresponding weighted and power inequalities}
Determine the valid exponents and sharp constants when the geometric mean
is weighted or the outer square is replaced by a positive power.

\begin{thebibliography}{9}
\bibitem{MF} M. S. Moslehian and X. Fu,
\emph{Squaring operator P\'olya--Szeg\H{o} and Diaz--Metcalf type inequalities},
Linear Algebra Appl. 491 (2016), 73--82.
\href{https://arxiv.org/abs/1501.02939}{arXiv:1501.02939v1},
Conjecture 2.16, p. 9.
\end{thebibliography}
\end{document}
